
\subsubsection{6.1 Key Findings and Implications}

\textbf{Finding 1: Scaling and sign-flip symmetry orbits are geometrically large in the Schürholt MNIST zoo --- the empirical motivation for canonicalization is established.}

Scaling orbits (mean cosine distance 0.32) and sign-flip orbits (mean cosine distance 1.07) are not geometric curiosities confined to pathological configurations --- they are universal properties of the 500-model Schürholt MNIST subset we study. Every oracle orbit pair exceeds the 0.05 significance threshold. This transforms the motivation for canonicalization from a theoretical argument to an empirical baseline: any method processing raw MLP weights for this zoo contends with within-orbit diameter of at least 0.32 (scaling) or 1.07 (sign-flip) cosine distance.

\textbf{Finding 2: NFT's symmetry handling is asymmetric --- non-invariant to scaling, measurably invariant to sign-flip by construction.}

For scaling orbits, NFT embeds functionally identical networks at different weight scales as significantly less similar than property-matched unrelated models (gap=+0.0238, CI entirely above zero). For sign-flip orbits, the opposite: NFT embeds orbit pairs as significantly more similar than cross-orbit property-matched pairs (gap=-0.00069, CI entirely below zero). This asymmetry was not built by design but emerges from NFT's row-level weight tokenization: sign-flip negates all entries of a weight row uniformly, leaving magnitude statistics unchanged, while scaling changes the magnitude structure in ways that affect NFT's attention mechanism.

This finding is actionable: for NFT-family encoders, scaling canonicalization addresses a measurable representational inefficiency. Sign-flip canonicalization, by contrast, may provide no benefit --- the encoder already treats sign-flip orbit members as similar. The recommendation is scaling-only canonicalization (Condition B) as the productive preprocessing step for NFT.

\textbf{Finding 3: The majority-sign sign-flip algorithm is structurally non-unique for even-d\_in architectures --- a previously uncharacterized limitation.}

The 85.6\% tie rate for d\_in=784 is a mathematical consequence of d\_in being even and of the approximately symmetric weight distributions in gradient-trained networks. This finding is not an implementation issue but a combinatorial property: any application of majority-sign canonicalization to 784-dimensional input MLPs will encounter this problem. The finding affects all weight space learning methods that apply sign-flip canonicalization as a preprocessing step for even-d\_in architectures. Three fixes are available: (1) use odd d\_in (e.g., pad MNIST inputs from 784 to 785); (2) use a secondary tie-breaking criterion based on weight magnitude; (3) use scaling-only canonicalization, which avoids the sign-flip problem entirely and may be sufficient given NFT's emergent sign-flip invariance.

\textbf{Finding 4: The E > D pattern (random normalization beats full canonicalization) is a cautionary result that cannot be confirmed at N=500.}

The consistent ordering E > D across all three properties and all 3 seeds is noteworthy but not statistically confirmable given CI width $\approx 0.6$. The most likely explanation is the non-unique sign-flip in Condition D (H-C1), which applied a tie-breaking convention to 85.6\% of models, potentially introducing structured noise. However, the N=400 underpowered training regime makes any specific attribution speculative. Until replicated at $N\geq 5,000$ with a corrected sign-flip implementation, practitioners should prefer Condition B (scaling-only) over Condition D.

\subsubsection{6.2 Limitations}

\textbf{Limitation 1: Zoo scale (N=500) provides insufficient statistical power for property prediction experiments.}

All property prediction experiments were conducted on N=500 models. At n=50 test samples, Spearman $\rho$ bootstrap CIs have width $\approx 0.6$ --- approximately $12\times$ wider than the effect size sought ($\Delta\rho \geq$ 0.05). No property prediction comparison in this paper has statistical power to distinguish conditions; all $\rho$ values are statistically indistinguishable from zero. This limitation does not invalidate the orbit characterization (H-E1) or the invariance probe (H-M1), which use relative within-orbit vs. cross-orbit comparisons that produce tight CIs (width $\approx 0.001)$ even at N=500. The full Schürholt zoo ($N\approx 50,000)$ would provide adequate power but was inaccessible via HuggingFace at experiment time.

\textbf{Limitation 2: Condition D in H-M3 tested a non-unique canonical transformation for 85.6\% of models.}

As described in Section 5.4, the majority-sign algorithm applied a +1 tie-break convention to the majority of zoo models. This means H-M3 Condition D did not test a true symmetry-canonical transformation. The P2 refutation (E > D) may reflect this artifact rather than a genuine failure of canonicalization. The proper comparison --- scaling-only (Condition B) versus Condition E at adequate N --- has not been performed.

\textbf{Limitation 3: Sign-flip functional equivalence in H-M1 is approximate, not exact.}

As described in Section 3.2, the simultaneous two-layer sign-flip is an approximate, not exact, functional equivalence under ReLU. Runtime functional equivalence verification on held-out inputs was not performed. The near-invariance result for sign-flip (gap=-0.0007) may partly reflect the approximate nature of the transform rather than NFT's architectural properties.

\textbf{Limitation 4: All findings derive from a single architecture and dataset.}

All experiments use a single zoo (MNIST), single architecture ($784\rightarrow 64\rightarrow 10$ MLP), and single dataset. The methodology (oracle orbit construction, within-orbit vs. cross-orbit probing, uniqueness audit) generalizes to other architectures. The specific numerical findings --- orbit diameters 0.32 and 1.07, NFT gap +0.024, tie rate 85.6\% --- are architecture-specific. Generalization to deeper networks, convolutional architectures, or non-ReLU activations is unknown.

\textbf{Limitation 5: NFT capacity as a binding constraint is not verified.}

The mechanistic argument assumes NFT's representational capacity is partially occupied by symmetry-induced variation. At N=500, NFT training is severely underpowered (val $\rho \approx$ 0), making it impossible to distinguish capacity constraint from data scarcity. With the full Schürholt zoo, a direct comparison of NFT val $\rho$ with and without canonicalization would isolate the capacity-constraint hypothesis.

\subsubsection{6.3 Future Work}

\textbf{Immediate:}

\begin{enumerate}
\item \textit{Full Schürholt zoo with scaling-only canonicalization (Condition B).} Acquires adequate statistical power ($N\geq 5,000)$ and avoids the sign-flip non-uniqueness issue. This single experiment provides a clean test of the core mechanism.
\end{enumerate}

\begin{enumerate}
\item \textit{Odd d\_in or magnitude-based tie-breaking for sign-flip canonicalization.} Padding MNIST inputs from 784 to 785 eliminates the tie problem and enables a clean test of full canonicalization (Condition D) against scaling-only (Condition B) to measure the marginal sign-flip contribution.
\end{enumerate}

\begin{enumerate}
\item \textit{Two-layer joint sign-flip functional equivalence verification.} Systematic runtime verification of functional equivalence for sign-flip oracle pairs used in H-M1, to confirm the sign-flip invariance result is not an artifact of an imprecise transform.
\end{enumerate}

\textbf{Medium-term:}

\begin{enumerate}
\item \textit{Architecture-specific symmetry audit for other encoders.} Apply the within-orbit vs. cross-orbit probing protocol to DWSNets \citep{Navon2023EquivariantArchitectures}, Universal Neural Functionals \citep{Kofinas2024UniversalNeuralFunctionals}, and Hyper-Representations [Schürholt et al., 2021] to characterize which symmetries each encoder naturally handles.
\end{enumerate}

\begin{enumerate}
\item \textit{Deeper networks (M > 2).} Sign-flip canonicalization for M > 2-layer MLPs is underdetermined by the majority-sign algorithm. Extending to M > 2 requires layer-by-layer greedy canonicalization or a more principled approach. The H-E1 and H-M1 infrastructure generalizes to M > 2 with minor modifications.
\end{enumerate}

